% AUTO-GENERATED BY Experiment pipeline. DO NOT EDIT BY HAND.
\section*{Khối Manuscript Tự Động (Results + Discussion ngắn)}

\subsection*{Results Block}
\noindent Trên 278 kích thước hiệu ứng từ 67 bài báo (90 thí nghiệm, giai đoạn 2020--2024), mô hình gộp 3 mức cho thấy Cộng hưởng mạnh âm ($g=-0.448$, 95\% CI [-0.664, -0.233], $p<0.001$), Tăng cường cho con người dương rõ rệt ($g=0.563$, 95\% CI [0.443, 0.682], $p<0.001$), và Tăng cường cho AI không có ý nghĩa thống kê ($g=0.058$, 95\% CI [-0.351, 0.467], $p=0.780$). Độ dị hợp cao ở cả ba đối sánh ($I^2$ lần lượt 97.8\%, 91.3\%, 99.4\%).

\subsection*{Bias Block}
\noindent Các kiểm định sai lệch công bố cho kết quả nhất quán: Egger có ý nghĩa ở 3/3 đối sánh, và trim-and-fill ước lượng số nghiên cứu thiếu dao động từ 11 đến 29. Điều này cho thấy cần diễn giải thận trọng các hiệu ứng thuận lợi trong tài liệu đã công bố.

\subsection*{Moderator + Industry Block}
\noindent Phân tích điều tiết cho thấy khác biệt theo bối cảnh rất lớn. Trong đối sánh Tăng cường cho AI, mức cao nhất xuất hiện ở Loại nhiệm vụ=Sáng tạo ($g=1.114$), trong khi mức thấp nhất ở Loại AI=AI mô phỏng ($g=-1.272$). Theo ngành, cùng đối sánh này dao động từ Truyền thông ($g=1.011$) đến Kinh doanh ($g=-0.732$), với $Q$-between $p<0.001$. Hàm ý là kết luận triển khai cần theo điều kiện nhiệm vụ và miền ứng dụng, thay vì một quy tắc chung cho mọi ngữ cảnh.

\noindent\textit{Nguồn tự động:} \texttt{Experiment/outputs/manifests/snapshot\_iter\_*.json}, \texttt{Experiment/outputs/reports/RESEARCH\_KEY\_FINDINGS\_VI.md}.
